\documentclass[10pt,a4paper]{amsart}
\usepackage[margin=19mm]{geometry}
\usepackage{amsmath,amssymb,mathtools}
\usepackage[scr=rsfs,cal=euler]{mathalfa}
\usepackage{xfrac}
\usepackage[unicode,hidelinks,bookmarks=false]{hyperref}
\newcommand{\ie}{i.e.\ }
\newcommand{\cf}{cf.\ }
\newcommand{\resp}{respectively\ }
\newcommand{\Subsection}{\subsection}
\providecommand{\nobreakdash}{\nobreak\hbox{-}}
\setlength{\emergencystretch}{2em}
\multlinegap0pt
\usepackage{bm}
\usepackage{bbm}
\usepackage{dsfont}
\usepackage{xspace}
\usepackage{cancel}
\usepackage{mathtools}
\usepackage{amsthm}
\makeatletter
\@ifundefined{theorem}{\newtheorem{theorem}{Theorem}}{}
\@ifundefined{lemma}{\newtheorem{lemma}[theorem]{Lemma}}{}
\makeatother
\newcommand{\dif}{\mathrm{d}}
\title[Initial boundary value problem for one-dimensional hyperboli]{Initial boundary value problem for one-dimensional hyperbolic compressible Navier-Stokes equations}
\author{Yuxi Hu, Yachun Li}
\date{Source: arXiv:2508.01634v1 (CC BY 4.0)}
\begin{document}
\maketitle

\noindent\emph{Source and licence.} Initial boundary value problem for one-dimensional hyperbolic compressible Navier-Stokes equations. Version arXiv:2508.01634v1 (CC BY 4.0), distributed under the Creative Commons Attribution 4.0 International licence, as recorded in the arXiv licence metadata for this exact version and shown on the abstract page https://arxiv.org/abs/2508.01634v1. Full rights notice: licenses/arxiv-2508.01634-CC-BY-4.0.txt.\par\smallskip

\noindent\textit{Editorial scope.} Counted unit: the Lemma whose source label is \texttt{Le3.10} in the article (source0.tex lines 597--607, source label \texttt{Le3.10}), delivered together with the article's own complete original proof of it (source0.tex lines 609--743, one \texttt{proof} environment of 135 lines). No mathematics is written, repaired or re-derived: statement and proof are the article's own text line for line. Required context retained uncounted: 2.1 (lines 199--205); 3.19 (lines 473--479); 3.23 (lines 536--542). Every definition, display and label of those lines is retained unchanged. Omitted from this package and not redistributed: 1--198, 206--472, 480--535, 543--596, 608--608, 744--900. Nothing inside a retained excerpt is omitted or altered: every retained body is delivered exactly as the article's lines are written, so no omission receipt of any kind accompanies this package. Citations: the retained excerpts carry no citation command at all, so no bibliography entry of the article is transcribed and no reference is redistributed. Reference adaptation: none was needed and none was made; every reference inside the delivered unit resolves inside it, so no unresolved reference or citation is produced. Printed slips kept literal: no printed word, symbol, hypothesis or number of the counted statement or of its proof was altered. Layout and class: the article's own class is \texttt{amsart} with packages including amsfonts, amsmath, CJK, amssymb, latexsym, bm, xcolor, mathrsfs, amsthm; the wrapper is the lane's pinned \texttt{amsart} editorial wrapper at 10pt on A4, which restates the theorem-like environments the retained text needs and adds the article's own macro definitions that the retained text needs (\texttt{\textbackslash dif}), with extra wrapper packages (bm, bbm, dsfont, xspace, cancel, mathtools). The delivered numbering is the unit's own. Assets: no retained span contains a figure environment, a graphics-inclusion command, a PDF-page inclusion, a table or a TikZ picture; no image file is copied into this package, the package declares no asset dependency and no image file is redistributed with it. Licence: version arXiv:2508.01634v1 of this article is distributed under the Creative Commons Attribution 4.0 International licence, as recorded in the arXiv licence metadata for this exact version and shown on the abstract page https://arxiv.org/abs/2508.01634v1. A full rights notice is delivered at licenses/arxiv-2508.01634-CC-BY-4.0.txt. Characters: every retained excerpt is pure ASCII, so the delivered engine, XeLaTeX with its default Latin Modern text font, reports no missing character.
\begin{align}\label{2.1}
\begin{cases}
v_t^\epsilon=u_x^\epsilon,\\
u_t^\epsilon+p(v^\epsilon)_x=S_x^\epsilon,\\
\tau (S_t^\epsilon+\epsilon b(x) S_x^\epsilon)+vS^\epsilon=\mu u_x^\epsilon, 
\end{cases}
\end{align}

\begin{align}\label{3.19}
\begin{cases}
v_{ttx}=u_{txx},\\
u_{ttx}+p(v)_{txx}=S_{txx},\\
\tau \left(\frac{S_t}{v}\right)_{tx}+\left(\tau \epsilon b(x)\frac{S_x}{v} \right)_{tx}+S_{tx}=\left(\frac{\mu}{v} u_x\right)_{tx}.
\end{cases}
\end{align}

\begin{align}\label{3.23}
\begin{cases}
v_{txx}=u_{xxx},\\
u_{txx}+p(v)_{xxx}=S_{xxx},\\
\tau \left(\frac{S_t}{v}\right)_{xx}+\tau \epsilon \left(\frac{b(x)S_x}{v}\right)_{xx}+S_{xx}=\left(\frac{\mu u_x}{v}\right)_{xx}.
\end{cases}
\end{align}

\begin{lemma}\label{Le3.10}
There exists a constant $C$ such that
\begin{align}
&\int_0^t \frac{1}{2} (u_{xx}^2(t,1)+u_{xx}^2(t,0)+ u_{tx}^2(t, 1)+u_{tx}^2(t, 0) ) \dif t \nonumber\\
 &\le C\int_0^1 (u_{tx}^2+u_{tt}^2+v_{tx}^2+\epsilon  v_{xx}^2 )\dif x
+C  \int_0^t\int_0^1 (u_{tx}^2+u_{xx}^2+u_{tt}^2 )\dif x \dif t \nonumber\\
&+C \epsilon \int_0^t\int_0^1 (v_{xx}^2+u_{xx}^2)\dif x \dif t +\frac{1}{8\mu} \int_0^t\int_0^1 S_{xx}^2 \dif x \dif t
+\frac{\tau \epsilon}{4\mu} \int_0^t (S_{xx}^2(t,1)+S_{xx}^2(t, 0)) \dif t +CE^\frac{1}{2}\int_0^t \mathcal D(s)\dif s .
\end{align}

\end{lemma}

\begin{proof}
 For simplicity, denote $\bar b(x):=-b(x)$ with $\bar b(0)=1, \bar b (1)=-1$. Multiplying the  equation $\eqref{3.23}_3$ by $\frac{v}{\mu} \bar b(x) u_{xx}$ and integrating the result, we get
\begin{align}
\int_0^1 \tau \left(\frac{S_t}{v}\right)_{xx} \frac{v}{\mu} \bar b(x) u_{xx} \dif x+\int_0^1 S_{xx}\frac{v}{\mu}\bar b(x) u_{xx} \dif x 
-\int_0^1 \left(\frac{\mu u_x}{v}\right)_{xx} \frac{v}{\mu} \bar b(x) u_{xx} \dif x \nonumber\\
=-\int_0^1\left( \tau\epsilon b(x) \frac{S_x}{v}\right)_{xx} \frac{v}{\mu} \bar b(x) u_{xx}\dif x . \label{3.26}
\end{align}
We estimate each term of the above equation, respectively. Firstly, we have
\begin{align*}
\int_0^1 \tau\left(\frac{S_t}{v}\right)_{xx} \frac{v}{\mu} \bar b(x) u_{xx} \dif x
=\int_0^1 \tau \left(\frac{S_{txx}}{v}-\frac{2v_x}{v^2} S_{tx}+\frac{2v_x^2}{v^3} S_t-\frac{v_{xx}}{v^2} S_t\right)\frac{v}{\mu} \bar b(x) u_{xx}\dif x\\
\ge \frac{\tau}{\mu} \frac{\dif}{\dif t}\int_0^1 S_{xx} \bar b(x) u_{xx}\dif x -\frac{\tau}{\mu}\int_0^1 (u_{tx}+p(v)_{xx}) \bar b(x) u_{txx}\dif x -C E^\frac{1}{2}(t) \mathcal D(t)
\end{align*}
where
\begin{align*}
-\frac{\tau}{\mu} \int_0^1 u_{tx} \bar b(x) u_{txx}\dif x=\frac{\tau}{\mu} (u_{tx}^2(t,1)+u_{tx}^2(t, 0)) -\frac{2\tau}{\mu}\int_0^1  u_{tx}^2\dif x 
\end{align*}
and
\begin{align*}
&-\frac{\tau}{\mu} \int_0^1 p(v)_{xx} \bar b(x) u_{txx}\dif x=-\frac{\dif}{\dif t} \int_0^1 \frac{\tau}{\mu}  p(v)_{xx}\bar b(x) u_{xx} \dif x+\frac{\tau}{\mu}\int_0^1 p(v)_{xxt} \bar b(x)u_{xx}\dif x\\
&\ge-\frac{\dif}{\dif t}\int_0^1 \frac{\tau}{\mu} p_{xx}(v) \bar b u_{xx}\dif x +\frac{\tau}{\mu} \int_0^1 p^\prime(v)u_{xx}^2\dif x +\frac{\tau}{\mu} p^\prime(v) u_{xx}^2 \bar b(x)\big|_0^1-CE^\frac{1}{2}(t)\mathcal D(t)\\
 &\ge-\frac{\dif}{\dif t}\int_0^1 \frac{\tau}{\mu} p_{xx}(v) \bar b u_{xx}\dif x +\frac{\tau}{\mu} \int_0^1 p^\prime(v)u_{xx}^2\dif x-CE^\frac{1}{2}(t)\mathcal D(t),
\end{align*} 
since $p^\prime(v)<0$. 

For the last term on the left-hand side of equation \eqref{3.26}, we get
\begin{align*}
&-\int_0^1 \left(\frac{\mu u_x}{v}\right)_{xx} \frac{v}{\mu} \bar b(x) u_{xx}\dif x\\
&=-\int_0^1\left(\frac{1}{v} u_{xxx}-\frac{2v_x}{v^2} u_{xx}+\frac{2v_x^2}{v^3} u_x-\frac{v_{xx}}{v^2}u_x\right)v \bar b(x)u_{xx}\dif x\\
&\ge -\frac{1}{2}\bar b(x) u_{xx}^2 \big|_0^1-\int_0^1 u_{xx}^2\dif x-CE^\frac{1}{2}(t)\mathcal D(t)\\
&\ge \frac{1}{2} (u_{xx}^2(t,1)+u_{xx}^2(t, 0))-\int_0^1 u_{xx}^2 \dif x-CE^\frac{1}{2}(t)\mathcal D(t).
\end{align*}



For the right-hand side of the equation \eqref{3.26}, we have
\begin{align*}
&\int_0^1 \tau \epsilon \left(\frac{b(x)S_x}{v}\right)_{xx} \frac{v}{\mu} \bar b(x)u_{xx}\dif x \\
&=\int_0^1 \tau \epsilon \left( \frac{b(x)}{v}S_{xxx}+\frac{4}{v} S_{xx}-\frac{2 b(x) v_x}{v^2}S_{xx}-\frac{4v_x}{v^2}S_x-\frac{b(x)v_{xx}}{v^2}S_x+\frac{2 b(x)v_x^2}{v^3} S_x\right)\frac{v}{\mu} \bar b(x)u_{xx}\dif x\\
&\le \frac{\tau \epsilon}{\mu} \int_0^1 b(x) S_{xxx} \bar b(x) u_{xx} \dif x +\frac{4\tau \epsilon}{\mu} \int_0^1 S_{xx} u_{xx} \dif x +CE^\frac{1}{2}(t)\mathcal D(t).
\end{align*}
In view of  equations $\eqref{2.1}_1$-$\eqref{2.1}_2$, we have
\begin{align*}
&\int_0^1 b(x) S_{xxx} \bar b(x) u_{xx}\dif x=-\int_0^1 b^2(x)(u_{txx}+p(v)_{xxx})u_{xx}\dif x\\
&\le -\frac{\dif}{\dif t}\int_0^1 \frac{1}{2}b^2(x) u_{xx}^2\dif x-\int_0^1 b^2(x)p^\prime(v) v_{xxx}u_{xx} \dif x+CE^\frac{1}{2}(t)\mathcal D(t)\\
&\le -\frac{\dif}{\dif t} \int_0^1 \frac{1}{2} b^2(x)u_{xx}^2 \dif x-\left[b^2(x)p^\prime(v) v_{xx}u_{xx} \right]\big|_0^1+\int_0^1 (b^2(x)p^\prime(v) u_{xx})_x v_{xx}\dif x +CE^\frac{1}{2}(t)\mathcal D(t)\\
&\le \frac{\dif}{\dif t} \int_0^1 (-\frac{1}{2} b^2 u_{xx}^2+\frac{1}{2} b^2(x)p^\prime(v) v_{xx}^2)\dif x+\eta[ v_{xx}^2(t,1)+v_{xx}^2(t,0)]+C(\eta) [u_{xx}^2(t,1)+u_{xx}^2(t,0)]\\
&\qquad\qquad\qquad\qquad\qquad\qquad\qquad\qquad\qquad\qquad+4\int_0^1 b(x)p^\prime(v) u_{xx}v_{xx}\dif x+CE^\frac{1}{2}(t)\mathcal D(t).
\end{align*}
where $\eta$ is sufficiently small and to be chosen later.
 
Combining the above result, we get
\begin{align*}
&\frac{1}{2} (u_{xx}^2(t,1)+u_{xx}^2(t,0))+\frac{\tau}{\mu} (u_{tx}^2 (t,1)+u_{tx}^2(t,0))\\
&\le \frac{\dif}{\dif t} \int_0^1 \left(-\frac{\tau}{\mu} \bar b(x) (S_{xx}-p_{xx}(v))u_{xx}\dif x+\frac{\tau \epsilon}{2\mu} b^2(x)(u_{xx}^2-p^\prime(v) v_{xx}^2)\right)\dif x\\
&+ \int_0^1\left( \frac{\tau}{\mu}(2 u_{tx}^2-p^\prime(v) u_{xx}^2 )+ \frac{v}{\mu} \bar b(x) S_{xx} u_{xx}+ u_{xx}^2+\frac{4\tau\epsilon}{\mu}  (S_{xx}+b(x)p^\prime(v) v_{xx}) u_{xx}\right) \dif x+C E^\frac{1}{2}(t)\mathcal D(t)\\
&+\frac{\tau\epsilon}{\mu} \eta[ v_{xx}^2(t,1)+v_{xx}^2(t,0)]+\frac{\tau\epsilon}{\mu} C(\eta) [u_{xx}^2(t,1)+u_{xx}^2(t,0)] .
\end{align*}
On the other hand, from the momentum equation $\eqref{2.1}_2$, we have
\begin{align*}
v_{xx}^2(t,1)+v_{xx}^2(t,0)\le C (S_{xx}^2(t,1)+S_{xx}^2(t,0)+u_{tx}^2(t,1)+u_{tx}^2(t,0))+C E^\frac{1}{2}(t)\mathcal D(t).
\end{align*}
Therefore, by choosing $\eta=\frac{1}{4C} $, $\epsilon$ sufficiently small and  integrating over $(0, t)$, we get
\begin{align}
&\int_0^t \frac{1}{2} (u_{xx}^2(t,1)+u_{xx}^2(t,0)) \dif t \nonumber\\
&\le \int_0^1 (v_{tx}^2+\epsilon  v_{xx}^2 +C u_{tx}^2)\dif x+C  \int_0^t\int_0^1 (u_{tx}^2+u_{xx}^2 )\dif x \dif t+CE^\frac{1}{2}\int_0^t \mathcal D(s)\dif s \nonumber\\
&+C \epsilon \int_0^t\int_0^1 v_{xx}^2\dif x \dif t +\frac{1}{8\mu} \int_0^t\int_0^1 S_{xx}^2 \dif x \dif t
+\frac{\tau \epsilon}{4\mu} \int_0^t (S_{xx}^2(t,1)+S_{xx}^2(t, 0)) \dif t . \label{3.28}
\end{align}

% Estimate of u_{tx}|_\Omega
Next, we we give estimates of $u_{tx}$ on the boundary. Multiplying  $\eqref{3.19}_3$ by $\frac{v}{\mu} \bar b(x) u_{tx}$ and integrating the result, we get
\begin{align}
\int_0^1 \tau \left(\frac{S_t}{v}\right)_{tx} \frac{v}{\mu} \bar b(x) u_{tx} \dif x+\int_0^1 S_{tx}\frac{v}{\mu}\bar b(x) u_{tx} \dif x 
-\int_0^1 \left(\frac{\mu u_x}{v}\right)_{tx} \frac{v}{\mu} \bar b(x) u_{tx} \dif x \nonumber\\
=-\int_0^1\left( \tau\epsilon b(x) \frac{S_x}{v}\right)_{tx} \frac{v}{\mu} \bar b(x) u_{tx}\dif x . \label{3.29}
\end{align}
We estimate each term of the above equation, respectively. Firstly, we have
\begin{align*}
\int_0^1 \tau\left(\frac{S_t}{v}\right)_{tx} \frac{v}{\mu} \bar b(x) u_{tx} \dif x
=\int_0^1 \tau \left(\frac{S_{ttx}}{v}-\frac{2v_x}{v^2} S_{tt}+\frac{2v_xv_t}{v^3} S_t-\frac{v_{tx}}{v^2} S_t\right)\frac{v}{\mu} \bar b(x) u_{tx}\dif x\\
\ge \frac{\tau}{\mu} \frac{\dif}{\dif t}\int_0^1 S_{tx} \bar b(x) u_{tx}\dif x -\frac{\tau}{\mu}\int_0^1 (u_{tt}+p(v)_{tx}) \bar b(x) u_{ttx}\dif x -C E^\frac{1}{2}(t) \mathcal D(t),
\end{align*}
where
\begin{align*}
-\frac{\tau}{\mu} \int_0^1 u_{tt} \bar b(x) u_{ttx}\dif x=-\frac{2\tau}{\mu}\int_0^1  u_{tt}^2\dif x
\end{align*}
and
\begin{align*}
&-\frac{\tau}{\mu} \int_0^1 p(v)_{tx} \bar b(x) u_{ttx}\dif x=-\frac{\dif}{\dif t} \int_0^1 \frac{\tau}{\mu}  p(v)_{tx}\bar b(x) u_{tx} \dif x+\frac{\tau}{\mu}\int_0^1 p(v)_{ttx} \bar b(x)u_{tx}\dif x\\
&\ge-\frac{\dif}{\dif t}\int_0^1 \frac{\tau}{\mu} p_{tx}(v) \bar b u_{tx}\dif x +\frac{\tau}{\mu} \int_0^1 p^\prime(v)u_{tx}^2\dif x +\frac{\tau}{\mu} p^\prime(v) u_{tx}^2 \bar b(x)\big|_0^1-CE^\frac{1}{2}(t)\mathcal D(t)\\
 &\ge-\frac{\dif}{\dif t}\int_0^1 \frac{\tau}{\mu} p_{tx}(v) \bar b u_{tx}\dif x +\frac{\tau}{\mu} \int_0^1 p^\prime(v)u_{tx}^2\dif x-CE^\frac{1}{2}(t)\mathcal D(t),
\end{align*} 
since $p^\prime(v)<0$. 

For the last term on the left-hand side of equation \eqref{3.28}, we get
\begin{align*}
&-\int_0^1 \left(\frac{\mu u_x}{v}\right)_{tx} \frac{v}{\mu} \bar b(x) u_{tx}\dif x\\
&=-\int_0^1\left(\frac{1}{v} u_{txx}-\frac{v_x}{v^2} u_{tx}-\frac{v_t}{v^2}u_{xx}+\frac{2v_x v_t}{v^3} u_x-\frac{v_{tx}}{v^2}u_x\right)v \bar b(x)u_{tx}\dif x\\
&\ge -\frac{1}{2}\bar b(x) u_{tx}^2 \big|_0^1-\int_0^1 u_{tx}^2\dif x-CE^\frac{1}{2}(t)\mathcal D(t)\\
&\ge \frac{1}{2} (u_{tx}^2(t,1)+u_{tx}^2(t, 0))-\int_0^1 u_{tx}^2 \dif x-CE^\frac{1}{2}(t)\mathcal D(t).
\end{align*}



For the right-hand side of the equation \eqref{3.28}, we have
\begin{align*}
&\int_0^1 \tau \epsilon \left(\frac{b(x)S_x}{v}\right)_{tx} \frac{v}{\mu} \bar b(x)u_{tx}\dif x \\
&=\int_0^1 \tau \epsilon \left( \frac{b(x)}{v}S_{txx}-\frac{ b(x) v_t}{v^2}S_{xx}-\frac{b(x) v_x}{v^2} S_{tx}-\frac{b(x)v_{tx}}{v^2}S_x+\frac{2 b(x)v_x v_t}{v^3} S_x\right)\frac{v}{\mu} \bar b(x)u_{tx}\dif x\\
&\le \frac{\tau \epsilon}{\mu} \int_0^1 b(x) S_{txx} \bar b(x) u_{tx} \dif x  +CE^\frac{1}{2}(t)\mathcal D(t).
\end{align*}
Using the  equations $\eqref{2.1}_1-\eqref{2.1}_2$, we have
\begin{align*}
&\int_0^1 b(x) S_{txx} \bar b(x) u_{tx}\dif x=-\int_0^1 b^2(x)(u_{ttx}+p(v)_{txx})u_{tx}\dif x\\
&\le -\frac{\dif}{\dif t}\int_0^1 \frac{1}{2}b^2(x) u_{tx}^2\dif x-\int_0^1 b^2(x)p^\prime(v) v_{txx}u_{tx} \dif x+CE^\frac{1}{2}(t)\mathcal D(t)\\
&\le -\frac{\dif}{\dif t} \int_0^1 \frac{1}{2} b^2(x)u_{tx}^2 \dif x-\left[b^2(x)p^\prime(v) u_{xx}u_{tx} \right]\big|_0^1+\int_0^1 (b^2(x)p^\prime(v) u_{tx})_x u_{xx}\dif x +CE^\frac{1}{2}(t)\mathcal D(t)\\
&\le \frac{\dif}{\dif t} \int_0^1 (-\frac{1}{2} b^2 u_{tx}^2+\frac{1}{2} b^2(x)p^\prime(v) u_{xx}^2)\dif x+C[ u_{xx}^2(t,1)+u_{xx}^2(t,0)+u_{tx}^2(t,1)+u_{tx}^2(t,0)]\\
&\qquad\qquad\qquad\qquad\qquad\qquad\qquad\qquad\qquad\qquad+4\int_0^1 b(x)p^\prime(v) u_{xx}u_{tx}\dif x+CE^\frac{1}{2}(t)\mathcal D(t).
\end{align*}

Combining the above result, we get
\begin{align*}
&\frac{1}{2} (u_{tx}^2(t,1)+u_{tx}^2(t,0))\le \frac{\dif}{\dif t} \int_0^1 \left(-\frac{\tau}{\mu} \bar b(x) (S_{tx}-p_{tx}(v))u_{tx}\dif x+\frac{\tau \epsilon}{2\mu} b^2(x)(u_{tx}^2-p^\prime(v) u_{xx}^2)\right)\dif x\\
&+C\epsilon\int_0^1 u_{xx}^2 \dif x+ \int_0^1\left( \frac{\tau}{\mu}(2 u_{tt}^2-p^\prime(v) u_{tx}^2 )+ \frac{v}{\mu} \bar b(x) S_{tx} u_{tx}+ u_{tx}^2\right) \dif x+C E^\frac{1}{2}(t)\mathcal D(t)\\
&+\frac{\tau\epsilon}{\mu} [ u_{xx}^2(t,1)+u_{xx}^2(t,0)].
\end{align*}
Integrating the above result over $(0, t)$, one obtain
\begin{align}
\int_0^t \frac{1}{2}( u_{tx}^2(t, 1)+u_{tx}^2(t, 0) )\dif t 
\le C(E_0+CE^\frac{1}{2}\int_0^t \mathcal D(s)\dif s+\int_0^1 (u_{tt}^2+u_{tx}^2+v_{tx}^2+\epsilon u_{xx}^2)\dif x \nonumber\\
+ \epsilon \int_0^t \int_0^1 u_{xx}^2 \dif x \dif t + \int_0^t\int_0^1 (u_{tt}^2+u_{tx}^2)\dif x \dif t 
+\epsilon \int_0^t ( u_{xx}^2(t,1)+u_{xx}^2(t,0)) \dif t) . \label{3.30}
\end{align}
Combining \eqref{3.28} and \eqref{3.30},  we get the desired result.
\end{proof}

\end{document}
